
\documentclass{article} %210 mm × 297 mm
\usepackage[utf8]{inputenc}
\usepackage[a4paper,left=1.5cm, right=1.5cm, top=2cm, bottom=2cm]{geometry}
\usepackage{graphicx}
\usepackage{systeme}          % for Solution 3
\usepackage{array}  
%\usepackage{blindtext}
\usepackage{multirow}
\usepackage[normalem]{ulem}
\usepackage{mhchem}
\usepackage{url}
\usepackage{subfigure} 
\usepackage{amsfonts,latexsym} % para tener disponibilidad de diversos simbolos
\usepackage{enumerate}
\usepackage{longdivision}
\usepackage{polynom}
%\usepackage{booktabs}
\usepackage{float}
%\usepackage{threeparttable}
\usepackage{array,colortbl}
%\usepackage{ifpdf}
%\usepackage{rotating}
\usepackage{stfloats}
\usepackage{url}
\usepackage{listings}
\usepackage{fancyhdr}
\usepackage{biblatex}
\usepackage{color}
\addbibresource{sources.bib}
%\usepackage{hyperref}
%\usepackage{wrapfig}
\usepackage{array}
%\usepackage{multirow}
%\usepackage{adjustbox}
%\usepackage{nccmath}
\usepackage[dvipsnames]{xcolor}
\usepackage{tcolorbox}
\newtcolorbox{mybox}{colback=white,
colframe=black}
\newcommand{\titulo}{Tutorium 10; Diskrete Mathematik}
\newcommand{\nombre}{Lil'Jana}
\newcommand{\FCE}{\textbf{WiSe 2024/2025}}

 

\usepackage{fancyhdr}
\pagestyle{fancy}
\fancyhead{}
\fancyfoot{}
\fancyhead[LO]{\small\nombre}
\fancyhead[RE]{\small\titulo}
\fancyhead[LO]{\small\FCE}
\fancyhead[RO]{\small 13.01.2025}
\fancyfoot[RO,LE] {\thepage}
\usepackage[onehalfspacing]{setspace}

\setcounter{page}{1} %Número de la página, la editorial lo cambiará
\begin{document}
\begin{tabular}{p{12cm}}
{{\Large\bf\titulo}}\\
\vspace{0.2cm}\\
 {\large\nombre}\\
 \vspace{0.1cm}
\end{tabular} 
\section*{Besprochene Lösungen}
\textbf{2a.} \\
monisch Grad 1. , $x +a $ für $a \in \mathbb{Z}/ p \mathbb{Z}$ \\
jedes solche Polynom muss irred. sein $x+a = f(x)g(x) \rightarrow 1 = deg f + deg g$  \\
OBdA $deg g = 0 \rightarrow g$ eine EInheit und $a+x$ ist irreduzibel. \\
$\Rightarrow$ $p = \#$ mon. irred. Grad! \\ \newline
\textbf{2b.}\\
monisch Grad 2, $x^2 +ax+b$, $a,b \in \mathbb{Z}/p \mathbb{Z}$, also p solche \\
redizibel, dann $(x+a\alpha)(x+\beta)$ für $\alpha, \beta \in \mathbb{Z}/ p \mathbb{Z}$
\begin{enumerate}
    \item Fall: $\alpha = \beta$: $p$ solche
    \item Fall: $\alpha \neq \beta$: $\frac{1}{2} p (p-1)$ solche
\end{enumerate}
$\Rightarrow p + \frac{1}{2}p(p-1) = \frac{1}{2}p(p+1)$ red Polynome \\
$\Rightarrow \# $ mon irred. Grad 2 $ = p^2 - \frac{1}{2}p(p+1) = \frac{1}{2}(p^2-p)$\\ 
\textbf{2c.} \\
Mon Polynome Grad 3 $x^3 + ax^2 + bx + c$, $a,b,c \in \mathbb{Z}/p \mathbb{Z}$, $p^3$ solche 
\begin{enumerate}
    \item Fall: $f(x) (x+\alpha)$ $f(x)$ mon. irred. Grad 2, $\alpha \in \mathbb{Z}/p \mathbb{Z} $ \\
    $\frac{1}{2}(p^2-p)\cdot p = \frac{1}{2} (p^3-p^2)$
    \item Fall: $(x+\alpha)(x+\beta)(x+\gamma) \quad \alpha, \beta, \gamma \in \mathbb{Z}/p \mathbb{Z}$ \\
    $(x+\alpha)^3$: $p$ solche \\
    $(x+\alpha)^3 (x+\beta), \alpha \neq \beta$: $p(p-1)$ solche \\
    $(x+\alpha)(x+\beta)(x+\gamma)$, $\alpha, \beta, \gamma$ paarweise verschieden $\frac{1}{6} p (p-1) (p-2)$ \\
    $\Rightarrow p+(p^2) + \frac{1}{6} (p^3-3p^2+3p) = \frac{1}{6} (p^3+3p^2+3p $ solche in Fall 2 
\end{enumerate}
Insgesamt gibt es $\frac{1}{6}(4p^3+2p) = \frac{1}{3} (2p^3+p) = \frac{1}{3}(p^3-p)$ \\ \newline
\textbf{3.} \\
C linear, dim k, Länge n, Minabstand $\delta$, Kontrollmatrix H  \\
\begin{tabular}{ll}
    linear: & $O \in C, \quad a,b \in C \Rightarrow a+b \in C$ \\
    dim k: &  entweder $\exists$ Basis mit k Elemente oder $|C| = 2^k$ \\
    Länge n: & jedes $c \in C$ hat Länge n, also besteht aus n Bits \\
    Minabst. $\delta$: & $\delta = \underset{min}{a \neq b \in C} \partial (a,b) \underbrace{=}_{falls \ linear} \underset{min}{a \neq 0 \in C} \underbrace{\omega(a)}_{\text{Gewicht = Anzahl der Einsen)}}$  \\
    Kontrollmatrix H: & hat Größe $(n-k) \times n$ \\
    & $C = \{x \mid H x^T = 0\}$
\end{tabular} \\ 
d:= max \{r | jede Teilmenge von den Spalten in H mit Größe r-1; ist linear unabhängig\} = min \{r | $\exists$ r linear abhängige Spalten in H\} \\
\textbf{a.} z.z. $\delta = d$  \\
Sei: \\
$H_i$ die i-te Spalte von H \\
$e_i = 0 ... 0 \underbrace{1}_{\text{i-te Pos.}} 0 ... 0 $ 
\begin{itemize}
    \item $H_i + ... + H_{ir} = 0 $ (also lin abhängig) $ \quad \Rightarrow $ Für $x = e_{i1} + ... + e_{i d}$ gilt: $Hx = 0$. ALso $x \in C$. ($d = \omega (x) \geq \delta)$ \\
    d minimal $\Rightarrow H_{i1} + ... + H_{i \delta} = 0$. \\
    Außerdem folgt aus d minimal noch: $\exists $ lin Komb. $a_i H_{i1} + ... + a_d H_{i d} = 0$ mit $a_j \in \{0,1\}$, nicht alle 0. \\
    Wenn $a_j = 0$ wäre --- (hier fehlt etwas)
    \item $\exists c \in C$ mit $\omega (c) = \delta $, schreibe $e_{i1} + ... + e_{i \delta}$ \\
    Dann gilt $Hc = 0$ Also $H_{i1} + ... + H_{i \delta} = 0$. \\
    $\Rightarrow H_{i1} + ... + H_{i \delta}$ sind lin. abh. $\Rightarrow d \leq \delta$ 
\end{itemize}
Rang H = max \{ r | $\exists$ r Spalten die lin. unabh. sind\} = n-k \\
$\delta $ = min \{ r | $\exists$ r Spalten, die lin abhängig sind\} \\
\textbf{b.} z.z. $\delta \leq n-k +1$ \\
Mögliche Beweise: 
\begin{itemize}
    \item Jede Teilmenge der Spalten der Größe $\delta-1$ ist linear unabhängig \\
$n-k = Rang H = max \{r | \exists r$ lin unabhängieg Spalten \} \\
Also $\delta-1 \leq n-k$
\item Entferne die letzten $\delta -1$ Einträge von jedem Codewort. Resultat ist der lin. Code C' der Länge $n-\delta +1$ mit gleicher Dimension k, da die Abb $C \longrightarrow C'$ ist eine Bijektion  \\
$\Rightarrow |C'| = 2^k \leq 2^{n-\delta+1}$, Also $k \leq n-\delta+1$
\end{itemize}
\textbf{4.} z.z. C alle Wörter der Länge n mit geradem Gewicht ist ein MDS Code \\
\textbf{linear}: \\
\begin{tabular}{ll}
    $O \in C$ & erfüllt  \\
    $a,b \in C $ & $\omega(a,b) = \omega(a) + \omega(b) \mod 2$ \\
    & also $\omega(a) , \omega(b)$ gerade $\Rightarrow \omega(a+b)$ gerade $\Rightarrow a+b \in C$ \\
    Alternativ:\\
    $\exists$ lin Abb. & $(\mathbb{Z}/ 2 \mathbb{Z})^n \rightarrow (\mathbb{Z}/ 2 \mathbb{Z})$ \\ 
    & $a \mapsto \omega(a) \mod 2$ \\
    & Ker = C
\end{tabular} \\
$\mathbf{\delta = 2}$ \\
$c = 1 1 0 ... 0$ in $C$, Gewicht 2 \\
$\delta = min \omega(a) $ gerade. ALso $\delta \neq 1 \Rightarrow \delta = 2$
\\ 
$\mathbf{k = n-1}$ \\
Bij. $C \overset{1-zu-1}{\leftrightarrow}$ \{Wörter mit ungeradem Gewicht\} \\
$C \mapsto c + e_1$ \\
$x \mapsto x+e_1 $ \\ 
$2^n = |Worter| = |C|+|\neg C| = 2|C|$ ALso $|C| = 2^{n-1}$ \\
\newline
Eine Kontrollmatrix ist nicht einheitlich und man muss eine Kontrollmatrix aufstellen können (KLausur)
\end{document}